% Options for packages loaded elsewhere
% Options for packages loaded elsewhere
\PassOptionsToPackage{unicode}{hyperref}
\PassOptionsToPackage{hyphens}{url}
\PassOptionsToPackage{dvipsnames,svgnames,x11names}{xcolor}
%
\documentclass[
  11pt,
]{article}
\usepackage{xcolor}
\usepackage[margin=1in]{geometry}
\usepackage{amsmath,amssymb}
\setcounter{secnumdepth}{-\maxdimen} % remove section numbering
\usepackage{iftex}
\ifPDFTeX
  \usepackage[T1]{fontenc}
  \usepackage[utf8]{inputenc}
  \usepackage{textcomp} % provide euro and other symbols
\else % if luatex or xetex
  \usepackage{unicode-math} % this also loads fontspec
  \defaultfontfeatures{Scale=MatchLowercase}
  \defaultfontfeatures[\rmfamily]{Ligatures=TeX,Scale=1}
\fi
\usepackage{lmodern}
\ifPDFTeX\else
  % xetex/luatex font selection
\fi
% Use upquote if available, for straight quotes in verbatim environments
\IfFileExists{upquote.sty}{\usepackage{upquote}}{}
\IfFileExists{microtype.sty}{% use microtype if available
  \usepackage[]{microtype}
  \UseMicrotypeSet[protrusion]{basicmath} % disable protrusion for tt fonts
}{}
\makeatletter
\@ifundefined{KOMAClassName}{% if non-KOMA class
  \IfFileExists{parskip.sty}{%
    \usepackage{parskip}
  }{% else
    \setlength{\parindent}{0pt}
    \setlength{\parskip}{6pt plus 2pt minus 1pt}}
}{% if KOMA class
  \KOMAoptions{parskip=half}}
\makeatother
% Make \paragraph and \subparagraph free-standing
\makeatletter
\ifx\paragraph\undefined\else
  \let\oldparagraph\paragraph
  \renewcommand{\paragraph}{
    \@ifstar
      \xxxParagraphStar
      \xxxParagraphNoStar
  }
  \newcommand{\xxxParagraphStar}[1]{\oldparagraph*{#1}\mbox{}}
  \newcommand{\xxxParagraphNoStar}[1]{\oldparagraph{#1}\mbox{}}
\fi
\ifx\subparagraph\undefined\else
  \let\oldsubparagraph\subparagraph
  \renewcommand{\subparagraph}{
    \@ifstar
      \xxxSubParagraphStar
      \xxxSubParagraphNoStar
  }
  \newcommand{\xxxSubParagraphStar}[1]{\oldsubparagraph*{#1}\mbox{}}
  \newcommand{\xxxSubParagraphNoStar}[1]{\oldsubparagraph{#1}\mbox{}}
\fi
\makeatother


\usepackage{longtable,booktabs,array}
\usepackage{calc} % for calculating minipage widths
% Correct order of tables after \paragraph or \subparagraph
\usepackage{etoolbox}
\makeatletter
\patchcmd\longtable{\par}{\if@noskipsec\mbox{}\fi\par}{}{}
\makeatother
% Allow footnotes in longtable head/foot
\IfFileExists{footnotehyper.sty}{\usepackage{footnotehyper}}{\usepackage{footnote}}
\makesavenoteenv{longtable}
\usepackage{graphicx}
\makeatletter
\newsavebox\pandoc@box
\newcommand*\pandocbounded[1]{% scales image to fit in text height/width
  \sbox\pandoc@box{#1}%
  \Gscale@div\@tempa{\textheight}{\dimexpr\ht\pandoc@box+\dp\pandoc@box\relax}%
  \Gscale@div\@tempb{\linewidth}{\wd\pandoc@box}%
  \ifdim\@tempb\p@<\@tempa\p@\let\@tempa\@tempb\fi% select the smaller of both
  \ifdim\@tempa\p@<\p@\scalebox{\@tempa}{\usebox\pandoc@box}%
  \else\usebox{\pandoc@box}%
  \fi%
}
% Set default figure placement to htbp
\def\fps@figure{htbp}
\makeatother





\setlength{\emergencystretch}{3em} % prevent overfull lines

\providecommand{\tightlist}{%
  \setlength{\itemsep}{0pt}\setlength{\parskip}{0pt}}



 


\usepackage{amsmath,amssymb,enumitem,booktabs}
\setlist{itemsep=2pt,topsep=3pt}
\usepackage{fancyhdr}\pagestyle{fancy}\fancyhf{}
\fancyhead[L]{\small DS701 Fall 2026}\fancyhead[R]{\small Homework 2}
\fancyfoot[C]{\small\thepage}
\renewcommand{\headrulewidth}{0.4pt}
\usepackage[breakable]{tcolorbox}
% --- answer-space macros (problems PDF only; no-ops in the solutions PDF) ---
% \answerspace{2in}  vertical room for handwritten work, sized from the solution
% \blank             a short fill-in rule; \blank[1.5in] for a wider one
% \answerpage        start the next problem on a fresh page (Gradescope page matching)
\newcommand{\answerspace}[1]{\par\vspace*{#1}}  % \vspace* so room is not dropped at a page break
\newcommand{\blank}[1][1.2in]{\rule{#1}{0.4pt}}
\newcommand{\answerpage}{\clearpage}
\newtcolorbox{solutionbox}{breakable,colback=blue!4,colframe=blue!45!black,title=Solution,fonttitle=\bfseries,left=4pt,right=4pt,top=3pt,bottom=3pt}
\makeatletter
\@ifpackageloaded{caption}{}{\usepackage{caption}}
\AtBeginDocument{%
\ifdefined\contentsname
  \renewcommand*\contentsname{Table of contents}
\else
  \newcommand\contentsname{Table of contents}
\fi
\ifdefined\listfigurename
  \renewcommand*\listfigurename{List of Figures}
\else
  \newcommand\listfigurename{List of Figures}
\fi
\ifdefined\listtablename
  \renewcommand*\listtablename{List of Tables}
\else
  \newcommand\listtablename{List of Tables}
\fi
\ifdefined\figurename
  \renewcommand*\figurename{Figure}
\else
  \newcommand\figurename{Figure}
\fi
\ifdefined\tablename
  \renewcommand*\tablename{Table}
\else
  \newcommand\tablename{Table}
\fi
}
\@ifpackageloaded{float}{}{\usepackage{float}}
\floatstyle{ruled}
\@ifundefined{c@chapter}{\newfloat{codelisting}{h}{lop}}{\newfloat{codelisting}{h}{lop}[chapter]}
\floatname{codelisting}{Listing}
\newcommand*\listoflistings{\listof{codelisting}{List of Listings}}
\makeatother
\makeatletter
\makeatother
\makeatletter
\@ifpackageloaded{caption}{}{\usepackage{caption}}
\@ifpackageloaded{subcaption}{}{\usepackage{subcaption}}
\makeatother
\usepackage{bookmark}
\IfFileExists{xurl.sty}{\usepackage{xurl}}{} % add URL line breaks if available
\urlstyle{same}
\hypersetup{
  pdftitle={DS701 Homework 2: Distances},
  colorlinks=true,
  linkcolor={blue},
  filecolor={Maroon},
  citecolor={Blue},
  urlcolor={Blue},
  pdfcreator={LaTeX via pandoc}}


\title{DS701 Homework 2: Distances}
\usepackage{etoolbox}
\makeatletter
\providecommand{\subtitle}[1]{% add subtitle to \maketitle
  \apptocmd{\@title}{\par {\large #1 \par}}{}{}
}
\makeatother
\subtitle{Released Mon Sep 14 · due Mon Sep 21, 11:59 pm on Gradescope}
\author{}
\date{}
\begin{document}
\maketitle


\begin{center}\LARGE\bfseries\color{blue!45!black} SOLUTIONS\end{center}
\renewcommand{\answerspace}[1]{}
\renewcommand{\answerpage}{\par\vspace{6pt}}

\subsection{Problem 1: The metric decides who is
close}\label{problem-1-the-metric-decides-who-is-close}

\emph{Practices: Minkowski distances, cosine similarity, Hamming and
Jaccard --- and noticing that ``closest'' depends on the choice.}

Three short documents are represented as word-count vectors over the
vocabulary (\emph{data}, \emph{model}, \emph{error}):
\[\mathbf{A} = (1, 1, 1),\qquad \mathbf{B} = (3, 3, 3),\qquad \mathbf{C} = (6, 1, 1).\]

\begin{enumerate}
\def\labelenumi{(\alph{enumi})}
\tightlist
\item
  For each of the three pairs (A,B), (A,C), (B,C) compute the Minkowski
  distance \(\Vert\mathbf{x}-\mathbf{y}\Vert_p\) for \(p = 1\)
  (Manhattan), \(p = 2\) (Euclidean) and \(p = \infty\) (Chebyshev) ---
  nine numbers, arranged as a \(3\times3\) table. Under each of the
  three distances, which pair is \textbf{closest}? Note where the answer
  changes.
\end{enumerate}

\begin{enumerate}
\def\labelenumi{(\alph{enumi})}
\setcounter{enumi}{1}
\tightlist
\item
  Compute the cosine similarity for the same three pairs (three
  decimals; \(\sqrt{3}\approx1.732\), \(\sqrt{27}\approx5.196\),
  \(\sqrt{38}\approx6.164\), \(\sqrt{114}\approx10.677\)). Which pair is
  most similar under cosine? Explain in one sentence why cosine gives A
  and B the value it does even though every \(\ell_p\) distance in (a)
  says they are some way apart, and say which of the two views is the
  right one if we want to know whether two documents are \emph{about the
  same thing}.
\end{enumerate}

\begin{enumerate}
\def\labelenumi{(\alph{enumi})}
\setcounter{enumi}{2}
\tightlist
\item
  Two pairs of shopping baskets, as sets of items:
\end{enumerate}

\begin{itemize}
\tightlist
\item
  \(P = \{\)bread, milk, eggs, butter, cheese, apples, rice,
  coffee\(\}\) and \(Q = \{\)bread, milk, eggs, butter, cheese, apples,
  rice, tea\(\}\);
\item
  \(R = \{\)bread, milk\(\}\) and \(S = \{\)bread, jam\(\}\).
\end{itemize}

For each pair compute the Hamming distance (the number of items in one
basket but not the other, i.e.~the size of the symmetric difference) and
the Jaccard similarity \(J_{\rm Sim} = |x\cap y|/|x\cup y|\). Which pair
does Hamming say is more alike, and which does Jaccard say? In one or
two sentences, say which measure you would trust for ``are these two
shoppers similar?'' and why.

\begin{solutionbox}

\begin{enumerate}
\def\labelenumi{(\alph{enumi})}
\tightlist
\item
  Difference vectors: \(\mathbf{A}-\mathbf{B} = (-2,-2,-2)\),
  \(\mathbf{A}-\mathbf{C} = (-5,0,0)\),
  \(\mathbf{B}-\mathbf{C} = (-3,2,2)\).
\end{enumerate}

\begin{longtable}[]{@{}cccc@{}}
\toprule\noalign{}
pair & \(p=1\) & \(p=2\) & \(p=\infty\) \\
\midrule\noalign{}
\endhead
\bottomrule\noalign{}
\endlastfoot
(A,B) & \(6\) & \(\sqrt{12}\approx3.46\) & \(2\) \\
(A,C) & \(5\) & \(\sqrt{25}=5\) & \(5\) \\
(B,C) & \(7\) & \(\sqrt{17}\approx4.12\) & \(3\) \\
\end{longtable}

Closest pair: \textbf{(A,C)} under Manhattan (5 \textless{} 6
\textless{} 7), but \textbf{(A,B)} under Euclidean (\(3.46\)) and under
Chebyshev (\(2\)). The ranking changes with \(p\): Manhattan adds up
three moderate differences for (A,B) and penalizes them more than the
single large difference of (A,C); Chebyshev sees only that single
largest difference.

\begin{enumerate}
\def\labelenumi{(\alph{enumi})}
\setcounter{enumi}{1}
\item
  \(\cos(\mathbf{A},\mathbf{B}) = \dfrac{3+3+3}{\sqrt3\,\sqrt{27}} = \dfrac{9}{\sqrt{81}} = 1.000\);
  \(\cos(\mathbf{A},\mathbf{C}) = \dfrac{6+1+1}{\sqrt3\,\sqrt{38}} = \dfrac{8}{\sqrt{114}} \approx \dfrac{8}{10.677} \approx 0.749\);
  \(\cos(\mathbf{B},\mathbf{C}) = \dfrac{18+3+3}{\sqrt{27}\,\sqrt{38}} = \dfrac{24}{\sqrt{1026}} \approx \dfrac{24}{32.03} \approx 0.749\)
  (same as (A,C), because \(\mathbf{B} = 3\mathbf{A}\) has the same
  direction as \(\mathbf{A}\)). Most similar: \textbf{(A,B)}, with
  cosine exactly 1: \(\mathbf{B} = 3\mathbf{A}\) points in the same
  direction, and cosine measures direction only --- it is unchanged by
  rescaling either vector --- while every \(\ell_p\) distance measures
  the length of the difference vector, which is nonzero. For ``about the
  same thing'' cosine is the right view: B is A with every word used
  three times as often (a document three times as long on the same
  topic), and topic is a matter of word \emph{proportions}, not counts.
\item
  \(P\cap Q\) has 7 items, \(P\cup Q\) has 9, symmetric difference
  \(\{\)coffee, tea\(\}\): Hamming \(=2\),
  \(J_{\rm Sim} = 7/9 \approx 0.778\). \(R\cap S = \{\)bread\(\}\),
  \(R\cup S = \{\)bread, milk, jam\(\}\), symmetric difference
  \(\{\)milk, jam\(\}\): Hamming \(=2\),
  \(J_{\rm Sim} = 1/3 \approx 0.333\). Hamming calls the two pairs
  \emph{equally} alike (2 and 2); Jaccard says (P,Q) is far more alike
  (\(0.78\) vs \(0.33\)). Trust Jaccard: P and Q agree on 7 of 9 items
  and differ on one swap, whereas R and S agree on 1 of 3 --- Jaccard
  normalizes the disagreement by how much there was to agree on, Hamming
  counts raw disagreements regardless of basket size.
\end{enumerate}

\emph{Common mistake:} in (a), forgetting the square root for \(p=2\)
(reporting 12, 25, 17), or taking \(\max\) of the raw coordinates rather
than of the absolute differences for \(p=\infty\).

\end{solutionbox}

\answerpage

\subsection{Problem 2: Is it a metric?}\label{problem-2-is-it-a-metric}

\emph{Practices: the metric axioms, especially the triangle inequality,
and why they matter downstream.}

Recall the four metric axioms: \(d(x,x)=0\); \(d(x,y)>0\) for
\(x\ne y\); \(d(x,y)=d(y,x)\); and the triangle inequality
\(d(x,y)\le d(x,z)+d(z,y)\) for all \(x,y,z\).

\begin{enumerate}
\def\labelenumi{(\alph{enumi})}
\tightlist
\item
  Let \(P=(0,0)\), \(Q=(3,4)\), \(R=(6,0)\) in \(\mathbb{R}^2\). Compute
  the three Euclidean distances \(d(P,Q)\), \(d(Q,R)\), \(d(P,R)\) and
  verify the triangle inequality for this triple in all three ways (each
  distance is at most the sum of the other two). Then state, in one
  line, when the triangle inequality holds with \textbf{equality} for
  Euclidean distance, and give a triple of points that achieves it.
\end{enumerate}

\begin{enumerate}
\def\labelenumi{(\alph{enumi})}
\setcounter{enumi}{1}
\tightlist
\item
  Cosine \emph{distance} is
  \(d_{\cos}(\mathbf{x},\mathbf{y}) = 1 - \cos(\mathbf{x},\mathbf{y})\).
  The lecture says it is \textbf{not} a metric. Prove it: give three
  explicit vectors \(\mathbf{x},\mathbf{z},\mathbf{y}\) in
  \(\mathbb{R}^2\) for which
  \(d_{\cos}(\mathbf{x},\mathbf{y}) > d_{\cos}(\mathbf{x},\mathbf{z}) + d_{\cos}(\mathbf{z},\mathbf{y})\),
  and show the numbers (use \(\cos 45^\circ = 1/\sqrt2 \approx 0.707\)
  if it helps). \emph{Hint:} think about angles \(0^\circ\),
  \(45^\circ\), \(90^\circ\).
\end{enumerate}

\begin{enumerate}
\def\labelenumi{(\alph{enumi})}
\setcounter{enumi}{2}
\tightlist
\item
  Cosine distance fails a \textbf{second} axiom as well. Which one, and
  what is a two-vector example? (One line.)
\end{enumerate}

\begin{enumerate}
\def\labelenumi{(\alph{enumi})}
\setcounter{enumi}{3}
\tightlist
\item
  In two or three sentences: why does it matter, for nearest-neighbor
  search or for clustering, whether the dissimilarity you use satisfies
  the triangle inequality? Concretely, what can you conclude about
  \(d(q,b)\) from knowing \(d(q,a)\) is small and \(d(a,b)\) is large
  when \(d\) is a metric --- and what goes wrong when it is not?
\end{enumerate}

\begin{solutionbox}

\begin{enumerate}
\def\labelenumi{(\alph{enumi})}
\item
  \(d(P,Q) = \sqrt{9+16} = 5\); \(d(Q,R) = \sqrt{9+16} = 5\);
  \(d(P,R) = \sqrt{36+0} = 6\). Checks: \(d(P,R) = 6 \le 5+5 = 10\);
  \(d(P,Q) = 5 \le 6+5 = 11\); \(d(Q,R) = 5 \le 5+6 = 11\). All hold
  (strictly). Equality holds exactly when the middle point lies
  \textbf{on the straight segment} between the other two (collinear,
  \(z\) between \(x\) and \(y\)), e.g.~\(P=(0,0)\), \(Z=(3,4)\),
  \(Y=(6,8)\): \(d(P,Y) = 10 = 5 + 5\).
\item
  Take \(\mathbf{x} = (1,0)\), \(\mathbf{z} = (1,1)\),
  \(\mathbf{y} = (0,1)\): angles \(0^\circ\), \(45^\circ\),
  \(90^\circ\). \(d_{\cos}(\mathbf{x},\mathbf{y}) = 1 - 0 = 1\);
  \(d_{\cos}(\mathbf{x},\mathbf{z}) = 1 - \dfrac{1}{1\cdot\sqrt2} = 1 - 0.707 = 0.293\),
  and by symmetry \(d_{\cos}(\mathbf{z},\mathbf{y}) = 0.293\). So
  \(d_{\cos}(\mathbf{x},\mathbf{z}) + d_{\cos}(\mathbf{z},\mathbf{y}) = 0.586 < 1 = d_{\cos}(\mathbf{x},\mathbf{y})\):
  going through \(\mathbf{z}\) is a ``shortcut'', which a metric
  forbids. One violating triple is a complete proof. (Remark: the
  \emph{angle} itself, \(\theta = \arccos(\cos)\), does satisfy the
  triangle inequality --- here \(90^\circ \le 45^\circ + 45^\circ\); it
  is \(1-\cos\) that bends the scale.)
\item
  Positivity: \(d(x,y)>0\) for \(x\ne y\) fails, since any two vectors
  in the same direction have \(d_{\cos} = 0\) --- e.g.~\((1,0)\) and
  \((3,0)\): \(\cos = 3/(1\cdot3) = 1\), so \(d_{\cos} = 0\) although
  the vectors differ.
\item
  With a metric, small \(d(q,a)\) and large \(d(a,b)\) force
  \(d(q,b) \ge d(a,b) - d(q,a)\) to be large: ``close to something far
  from \(b\)'' means ``far from \(b\)'', so a search can skip \(b\)
  without computing \(d(q,b)\), and a cluster of mutually close points
  cannot contain a point far from all its members. Without the triangle
  inequality none of that follows --- \(q\) can be close to \(a\), \(a\)
  close to \(b\), and \(q\) far from \(b\) --- so pruning is unsafe and
  ``clusters'' of near neighbors need not be internally coherent.
  (Cosine distance is still used constantly; the point is to know which
  guarantees you have given up.)
\end{enumerate}

\emph{Common mistake:} trying to prove (b) with vectors like \((1,0)\)
and \((2,0)\) --- those have \(d_{\cos}=0\) and show positivity failing
(part (c)), not the triangle inequality.

\end{solutionbox}

\answerpage

\subsection{Problem 3: Same points, different units, different
neighbor}\label{problem-3-same-points-different-units-different-neighbor}

\emph{Practices: why raw Euclidean distance is dominated by the
largest-range feature, and what standardizing does.}

A bank describes customers by two features: annual income (in dollars)
and age (in years). A query customer \(q\) and four customers on file:

\begin{longtable}[]{@{}crr@{}}
\toprule\noalign{}
& income & age \\
\midrule\noalign{}
\endhead
\bottomrule\noalign{}
\endlastfoot
\(q\) & 50,000 & 30 \\
\(P_1\) & 52,000 & 60 \\
\(P_2\) & 60,000 & 32 \\
\(P_3\) & 45,000 & 28 \\
\(P_4\) & 51,000 & 45 \\
\end{longtable}

To \textbf{standardize} a feature, replace each value \(v\) by its
\(z\)-score \((v-\mu)/\sigma\), where \(\mu\) and \(\sigma\) are the
feature's mean and standard deviation. Use these values, taken from the
bank's full customer base: income \(\mu = 50{,}000\),
\(\sigma = 10{,}000\); age \(\mu = 40\), \(\sigma = 10\).

\begin{enumerate}
\def\labelenumi{(\alph{enumi})}
\tightlist
\item
  Using raw Euclidean distance, which of \(P_1,\dots,P_4\) is the
  nearest neighbor of \(q\)? (Squared distances are enough to rank;
  report them.) In one line: how much does the age difference contribute
  to any of these distances?
\end{enumerate}

\begin{enumerate}
\def\labelenumi{(\alph{enumi})}
\setcounter{enumi}{1}
\tightlist
\item
  Standardize all five rows with the given \(\mu,\sigma\) (write the
  \(z\)-scores), recompute the four squared Euclidean distances from
  \(q\), and give the nearest neighbor now. Does it change?
\end{enumerate}

\begin{enumerate}
\def\labelenumi{(\alph{enumi})}
\setcounter{enumi}{2}
\tightlist
\item
  Redo (a) with income measured in \emph{thousands} of dollars (50, 52,
  60, 45, 51) and age unchanged. Which neighbor wins? Then answer: which
  of the three answers should a data scientist trust as ``the customer
  most like \(q\)'', and why --- what is wrong with the other two?
\end{enumerate}

\begin{enumerate}
\def\labelenumi{(\alph{enumi})}
\setcounter{enumi}{3}
\tightlist
\item
  Standardizing is not always right. Give one situation where you would
  \textbf{not} standardize before computing distances, and one sentence
  on why (think about features that are already in the same meaningful
  units, or a feature whose spread \emph{should} count).
\end{enumerate}

\begin{solutionbox}

\begin{enumerate}
\def\labelenumi{(\alph{enumi})}
\item
  Squared raw distances from \(q\): \(P_1\):
  \(2000^2 + 30^2 = 4{,}000{,}000 + 900 = 4{,}000{,}900\); \(P_2\):
  \(10{,}000^2 + 2^2 = 100{,}000{,}004\); \(P_3\):
  \(5000^2 + 2^2 = 25{,}000{,}004\); \(P_4\):
  \(1000^2 + 15^2 = 1{,}000{,}225\). Nearest neighbor: \textbf{\(P_4\)}
  (\(d\approx1000.1\)). Age contributes at most \(900\) out of millions
  --- well under \(0.1\%\); the ranking is decided by income alone.
\item
  \(z\)-scores
  \(\big((\text{income}-50{,}000)/10{,}000,\ (\text{age}-40)/10\big)\):
  \(q = (0,\,-1)\), \(P_1 = (0.2,\,2)\), \(P_2 = (1,\,-0.8)\),
  \(P_3 = (-0.5,\,-1.2)\), \(P_4 = (0.1,\,0.5)\). Squared distances from
  \(q\): \(P_1\): \(0.04 + 9 = 9.04\); \(P_2\): \(1 + 0.04 = 1.04\);
  \(P_3\): \(0.25 + 0.04 = 0.29\); \(P_4\): \(0.01 + 2.25 = 2.26\).
  Nearest neighbor: \textbf{\(P_3\)} (\(d = \sqrt{0.29}\approx0.54\)).
  Yes --- it flips: \(P_4\), the raw winner, is now third, because its
  15-year age gap (1.5 SD) finally counts.
\item
  Income in thousands: \(P_1\): \(2^2 + 30^2 = 904\); \(P_2\):
  \(10^2 + 2^2 = 104\); \(P_3\): \(5^2+2^2 = 29\); \(P_4\):
  \(1^2 + 15^2 = 226\). Nearest: \textbf{\(P_3\)} --- a \emph{third}
  computation on the same customers, and now age dominates instead of
  income. Trust (b). The raw answers in (a) and (c) depend on an
  arbitrary bookkeeping choice --- dollars vs thousands of dollars ---
  that says nothing about customers; each lets whichever feature happens
  to have the bigger numbers decide. Standardizing puts both features on
  the same scale (SD units), so a 1-SD income gap and a 1-SD age gap
  count equally and the answer no longer changes when someone re-enters
  the spreadsheet in different units.
\item
  Any of: (i) features already in the \emph{same} meaningful units where
  spread is information --- e.g.~the pixel intensities of an image, or
  coordinates in meters, or the same measurement at 12 months:
  standardizing would inflate a nearly-constant, mostly-noise feature to
  equal importance with a genuinely varying one; (ii) one-hot/binary
  indicator features, whose 0/1 scale is already comparable; (iii) when
  the domain says one feature \emph{should} dominate. The rule:
  standardize when the units are arbitrary and incomparable; don't when
  the units are shared and the spread is the signal.
\end{enumerate}

\emph{Common mistake:} standardizing with the mean and SD of only the
five rows in the table when the problem gives population values ---
either is defensible only if you say which you used; here the given
values are the ones to use.

\end{solutionbox}




\end{document}
